\chapter{Android via the NDK}
\label{ch:android-ndk}

CNA's README states Android is a real, verified target --- and this chapter's job is to
describe precisely what that verification actually covers, because, as with the previous
chapter's account of the Web target, the honest picture has sharper edges than the one-line
summary suggests.

\section{Two separate build paths}

Android support in this project is not one build path but two, serving different purposes.
The first is a plain CMake cross-compile using the NDK's own toolchain file directly, with no
Gradle involved at all --- producing just the \texttt{CNA} static library, not the test suite,
since GoogleTest has never been configured for the NDK toolchain in this project's history.
Confirming that platform-conditional Android code paths (Chapter~\ref{ch:devices-sensors}'s
Android-specific \cnaclass{Compass}/\cnaclass{Motion} backends, for instance) actually
\emph{compile} under this path uses the NDK's own \texttt{llvm-nm} tool directly, since a
host Linux \texttt{nm} cannot read ARM64 object files.

The second path is a full Gradle-based Android application project, generated from SDL's own
project-creation template, then rewired so its native build step adds CNA's own root CMake
project as a subdirectory rather than building a second, separately vendored copy of SDL ---
a real, deliberate deduplication, since the template's own vendored SDL copy was roughly
34~MB before being removed in favor of reusing CNA's existing one.

\section{A real regression, caught honestly rather than left uncorrected}

This chapter's most important fact is a documented regression, and it is worth stating
plainly rather than glossing over: an earlier verification pass confirmed the plain
NDK cross-compile succeeded end to end. A later verification pass, re-running the identical
command, found it now fails --- and, tellingly, the failure is not in any CNA code at all,
but in two specific files inside the sibling \texttt{sharp-runtime} project
(Chapter~\ref{ch:sharp-runtime-overview}), before CNA's own target even begins compiling: a
strict-warnings build catching an unused private field on one platform/compiler
configuration that a desktop build does not, and a use of
\texttt{std::chrono::clock\_cast} that the NDK's bundled C\texttt{++} standard library does
not yet implement. Because the build never reaches CNA's own \texttt{SdlGraphicsBackend.cpp}
at all under this regression, that file's own Android buildability is honestly recorded as
currently \emph{unverified}, not confirmed working --- a direct, current example of exactly
the kind of documentation staleness Chapter~\ref{ch:easygl-backend} already flagged for
EasyGL's own bug list: a real success recorded at one point in time, superseded by a later,
less favorable finding, and this book reports the more recent one rather than the more
flattering one.

\section{What the Gradle path actually proves, concretely}

The one real Android application in this project is not a graphics demo at all --- it is a
devices/sensors demo, and no Android graphics demo project exists in the repository. Building
it surfaced and fixed several genuine, specific problems: a stale cached SDL-prebuilt path
setting; a missing \texttt{PUBLIC android} link dependency needed for the NDK's native sensor
manager and looper APIs; and a real launch failure (\texttt{Couldn't find function SDL\_main
in library libmain.so}) traced to a missing header whose inclusion is what actually redirects
a game's \texttt{main()} to the symbol SDL's own Java activity code looks up via
\texttt{dlsym()} at launch.

\section{What was actually observed running on the emulator}

The claim of ``verified on a real x86\_64 emulator'' resolves to a specific configuration: a
particular Android Virtual Device profile, run with software rendering and no hardware audio,
after confirming that hardware-accelerated virtualization was genuinely available inside this
project's own development container (an earlier attempt, in a container without that
acceleration available, failed to boot the emulator at all --- itself an honestly recorded,
environment-specific failure, not a CNA bug). Once booted, the devices demo was installed and
launched, and its liveness was confirmed three separate ways: a live process ID, a real
screen capture matching the demo's own expected layout (not a blank or crashed screen), and ---
most tellingly --- synthetic sensor values injected directly through the emulator's own
console (a specific acceleration vector, a specific magnetic-field vector) producing a real,
visible on-screen response, live evidence that SDL3's Android sensor event pipeline delivers
real events end to end, not merely that the app launched without crashing.

Two honest limits on this result are worth stating directly. First, this is evidence for the
devices/sensors demo specifically --- \texttt{SDL\_RENDERER} (Android's own default graphics
backend, since Android's own CMake system name is \texttt{"Android"}, not \texttt{"Linux"},
so it never falls into \texttt{EASYGL}'s Linux-default branch) has never actually been run on
this or any emulator or device, and its buildability is separately unverified for the reason
already given above. Second, the emulator's own virtual sensor set has no rotation-vector
sensor, so the Android-specific \cnaclass{Compass}/\cnaclass{Motion} backends
(Chapter~\ref{ch:devices-sensors}) were confirmed only to compile, link, and launch
successfully --- never actually exercised with injected sensor values the way
\cnaclass{Accelerometer}/\cnaclass{Gyroscope} were. No physical Android device has been used
for anything in this project's history at all; every claim in this chapter concerns the
emulator specifically.
